% GENERATED FILE. Edit canonical lessons, then run npm run generate:books.
% canonical-source-hash: fnv1a64:62da4ca095f07354
% canonical-lessons: JA-C338-muda, JA-C338-herasu, JA-C338-fuyasu, JA-C338-moyasu, JA-C338-shigen

\chapter{Less and More}
\label{ch:ja-less-and-more}
\hlchaptermodality{338}

\begin{chapteropening}
\textbf{By the end of this chapter:} \emph{I can say what we should cut down on and what we should increase.}

Complete the last lesson of chapter 338: I can say what we should cut down on and what we should increase.
\end{chapteropening}

\section[\texorpdfstring{\ja{むだ} (muda)}{むだ (muda)}]{\ja{むだ} (muda) — useless; a waste}

\label{lesson:JA-C338-muda}



\begin{quote}
Before the new one: say the Japanese for petroleum, then the Japanese for too good to waste.
\end{quote}

\subsection*{You'll want to know: \ja{むだ}}
\textbf{\ja{むだ}} — \emph{muda} — \textquotedblleft{}useless; a waste\textquotedblright{}.

\begin{quote}
\textbf{\ja{ひるに} \ja{でんきを} \ja{つかったら、むだです。}} — \emph{hiru ni denki o tsukattara, muda desu} — \textquotedblleft{}Using the lights in the daytime is a waste.\textquotedblright{}
\end{quote}

A na-adjective: \ja{な} before a noun.

\begin{cousinweb}
\textbf{\ja{むだ}} (\emph{muda}) is written as Sino-Japanese (kango): two Chinese characters read with their on'yomi, readings borrowed from Chinese, \emph{mu}, \textquotedblleft{}nothing\textquotedblright{} (Middle Chinese \emph{mju}) and \emph{da}, \textquotedblleft{}burdensome\textquotedblright{}. JMdict marks that spelling as ateji, characters chosen for their sound, so whether the word itself came from Chinese is not settled (JMdict, KANJIDIC2; Middle Chinese in Baxter's transcription).
\end{cousinweb}

\begin{grammarlens}[title={What you've built}]
One more word for giving reasons and opinions.
\end{grammarlens}

\subsection*{Your turn}
\begin{itemize}
\raggedright
  \item \emph{Say it:} \emph{muda}
  \item \emph{Say it:} \emph{muda}, once more
  \item \emph{Say it:} \emph{hiru ni denki o tsukattara, muda desu} — Using the lights in the daytime is a waste.
  \item \emph{Recall:} say the Japanese for petroleum, then the Japanese for too good to waste, then say \emph{muda} again
\end{itemize}

\begin{tcolorbox}[breakable,colback=teal!4,colframe=teal!35!black,title={Before you move on}]
What does \ja{むだ} mean? (\textquotedblleft{}Useless; a waste\textquotedblright{}.) Say it once more.
\end{tcolorbox}
\section[\texorpdfstring{\ja{へらす} (herasu)}{へらす (herasu)}]{\ja{へらす} (herasu) — to reduce; to cut down}

\label{lesson:JA-C338-herasu}



\begin{quote}
Before the new one: say the Japanese for too good to waste, then the Japanese for useless.
\end{quote}

\subsection*{You'll want to know: \ja{へらす}}
\textbf{\ja{へらす}} — \emph{herasu} — \textquotedblleft{}to reduce; to cut down\textquotedblright{}.

\begin{quote}
\textbf{\ja{かんきょうの} \ja{ために、ごみを} \ja{へらした} \ja{ほうが} \ja{いいです。}} — \emph{kankyō no tame ni, gomi o herashita hō ga ii desu} — \textquotedblleft{}For the environment, we should reduce our rubbish.\textquotedblright{}
\end{quote}

Polite \textbf{\ja{へらします}}, \emph{herashimasu}; a godan verb. Past (polite): \textbf{\ja{まちが} \ja{ごみを} \ja{へらしました。}}, \emph{machi ga gomi o herashimashita}.

\begin{cousinweb}
\textbf{\ja{へらす}} (\emph{herasu}) is a native Japanese word (wago), inherited rather than borrowed from Chinese, and its older origin is not settled. Its kanji was attached to it later as a kun'yomi, a native reading given to a Chinese character (JMdict, KANJIDIC2).
\end{cousinweb}

\begin{grammarlens}[title={What you've built}]
One more word for giving reasons and opinions.
\end{grammarlens}

\subsection*{Your turn}
\begin{itemize}
\raggedright
  \item \emph{Say it:} \emph{herasu}
  \item \emph{Say it:} \emph{herasu}, once more
  \item \emph{Say it:} \emph{kankyō no tame ni, gomi o herashita hō ga ii desu} — For the environment, we should reduce our rubbish.
  \item \emph{Recall:} say the Japanese for too good to waste, then the Japanese for useless, then say \emph{herasu} again
\end{itemize}

\begin{tcolorbox}[breakable,colback=teal!4,colframe=teal!35!black,title={Before you move on}]
What does \ja{へらす} mean? (\textquotedblleft{}To reduce; to cut down\textquotedblright{}.) Say it once more.
\end{tcolorbox}
\section[\texorpdfstring{\ja{ふやす} (fuyasu)}{ふやす (fuyasu)}]{\ja{ふやす} (fuyasu) — to increase something; to add to}

\label{lesson:JA-C338-fuyasu}



\begin{quote}
Before the new one: say the Japanese for useless, then the Japanese for to reduce.
\end{quote}

\subsection*{You'll want to know: \ja{ふやす}}
\textbf{\ja{ふやす}} — \emph{fuyasu} — \textquotedblleft{}to increase something; to add to\textquotedblright{}.

\begin{quote}
\textbf{\ja{けんこうの} \ja{ために、やさいを} \ja{ふやした} \ja{ほうが} \ja{いいです。}} — \emph{kenkō no tame ni, yasai o fuyashita hō ga ii desu} — \textquotedblleft{}For your health, you should eat more vegetables.\textquotedblright{}
\end{quote}

Polite \textbf{\ja{ふやします}}, \emph{fuyashimasu}; a godan verb. Past (polite): \textbf{\ja{ちちが} \ja{やさいを} \ja{ふやしました。}}, \emph{chichi ga yasai o fuyashimashita}.

\begin{cousinweb}
\textbf{\ja{ふやす}} (\emph{fuyasu}) is a native Japanese word (wago), inherited rather than borrowed from Chinese, and its older origin is not settled. Its kanji was attached to it later as a kun'yomi, a native reading given to a Chinese character (JMdict, KANJIDIC2).
\end{cousinweb}

\begin{grammarlens}[title={What you've built}]
One more word for giving reasons and opinions.
\end{grammarlens}

\subsection*{Your turn}
\begin{itemize}
\raggedright
  \item \emph{Say it:} \emph{fuyasu}
  \item \emph{Say it:} \emph{fuyasu}, once more
  \item \emph{Say it:} \emph{kenkō no tame ni, yasai o fuyashita hō ga ii desu} — For your health, you should eat more vegetables.
  \item \emph{Recall:} say the Japanese for useless, then the Japanese for to reduce, then say \emph{fuyasu} again
\end{itemize}

\begin{tcolorbox}[breakable,colback=teal!4,colframe=teal!35!black,title={Before you move on}]
What does \ja{ふやす} mean? (\textquotedblleft{}To increase something; to add to\textquotedblright{}.) Say it once more.
\end{tcolorbox}
\section[\texorpdfstring{\ja{もやす} (moyasu)}{もやす (moyasu)}]{\ja{もやす} (moyasu) — to burn something}

\label{lesson:JA-C338-moyasu}



\begin{quote}
Before the new one: say the Japanese for to reduce, then the Japanese for to increase something.
\end{quote}

\subsection*{You'll want to know: \ja{もやす}}
\textbf{\ja{もやす}} — \emph{moyasu} — \textquotedblleft{}to burn something\textquotedblright{}.

\begin{quote}
\textbf{\ja{ごみを} \ja{もやしたら、くうきが} \ja{よごれます。}} — \emph{gomi o moyashitara, kūki ga yogoremasu} — \textquotedblleft{}Burning rubbish makes the air dirty.\textquotedblright{}
\end{quote}

Polite \textbf{\ja{もやします}}, \emph{moyashimasu}; a godan verb. Past (polite): \textbf{\ja{おとこが} \ja{ごみを} \ja{もやしました。}}, \emph{otoko ga gomi o moyashimashita}.

\begin{cousinweb}
\textbf{\ja{もやす}} (\emph{moyasu}) is a native Japanese word (wago), inherited rather than borrowed from Chinese, and its older origin is not settled. Its kanji was attached to it later as a kun'yomi, a native reading given to a Chinese character (JMdict, KANJIDIC2).
\end{cousinweb}

\begin{grammarlens}[title={What you've built}]
One more word for giving reasons and opinions.
\end{grammarlens}

\subsection*{Your turn}
\begin{itemize}
\raggedright
  \item \emph{Say it:} \emph{moyasu}
  \item \emph{Say it:} \emph{moyasu}, once more
  \item \emph{Say it:} \emph{gomi o moyashitara, kūki ga yogoremasu} — Burning rubbish makes the air dirty.
  \item \emph{Recall:} say the Japanese for to reduce, then the Japanese for to increase something, then say \emph{moyasu} again
\end{itemize}

\begin{tcolorbox}[breakable,colback=teal!4,colframe=teal!35!black,title={Before you move on}]
What does \ja{もやす} mean? (\textquotedblleft{}To burn something\textquotedblright{}.) Say it once more.
\end{tcolorbox}
\section[\texorpdfstring{\ja{しげん} (shigen)}{しげん (shigen)}]{\ja{しげん} (shigen) — resources}

\label{lesson:JA-C338-shigen}



\begin{quote}
Before the new one: say the Japanese for to increase something, then the Japanese for to burn something.
\end{quote}

\subsection*{You'll want to know: \ja{しげん}}
\textbf{\ja{しげん}} — \emph{shigen} — \textquotedblleft{}resources\textquotedblright{}.

\begin{quote}
\textbf{\ja{しげんが} \ja{すくないので、せつやくした} \ja{ほうが} \ja{いいです。}} — \emph{shigen ga sukunai node, setsuyaku shita hō ga ii desu} — \textquotedblleft{}Resources are scarce, so we should use them carefully.\textquotedblright{}
\end{quote}

\begin{cousinweb}
\textbf{\ja{しげん}} (\emph{shigen}) is Sino-Japanese (kango): two Chinese characters read with their on'yomi, readings borrowed from Chinese. They are \emph{shi}, \textquotedblleft{}resources\textquotedblright{} (Middle Chinese \emph{tsij}) and \emph{gen}, \textquotedblleft{}source\textquotedblright{} (KANJIDIC2; Middle Chinese in Baxter's transcription).
\end{cousinweb}

\begin{grammarlens}[title={What you've built}]
One more word for giving reasons and opinions.
\end{grammarlens}

\subsection*{Your turn}
\begin{itemize}
\raggedright
  \item \emph{Say it:} \emph{shigen}
  \item \emph{Say it:} \emph{shigen}, once more
  \item \emph{Say it:} \emph{shigen ga sukunai node, setsuyaku shita hō ga ii desu} — Resources are scarce, so we should use them carefully.
  \item \emph{Recall:} say the Japanese for to increase something, then the Japanese for to burn something, then say \emph{shigen} again
\end{itemize}

\begin{tcolorbox}[breakable,colback=teal!4,colframe=teal!35!black,title={Before you move on}]
What does \ja{しげん} mean? (\textquotedblleft{}Resources\textquotedblright{}.) Say it once more.
\end{tcolorbox}
